\documentclass [a4paper,12pt]{paper}
%DIF LATEXDIFF DIFFERENCE FILE
%DIF DEL lecture3.tex             Wed Oct 29 04:34:40 2025
%DIF ADD diffs/lecture3_new.tex   Fri Sep 25 16:34:48 2026

\usepackage{mycommands}

\title{Lecture 14: Seismic sources}
\author{David Al-Attar \\ Michaelmas Term, \the\year}
%DIF PREAMBLE EXTENSION ADDED BY LATEXDIFF
%DIF UNDERLINE PREAMBLE %DIF PREAMBLE
\RequirePackage[normalem]{ulem} %DIF PREAMBLE
\RequirePackage{color}\definecolor{RED}{rgb}{1,0,0}\definecolor{BLUE}{rgb}{0,0,1} %DIF PREAMBLE
\providecommand{\DIFadd}[1]{{\protect\color{blue}\uwave{#1}}} %DIF PREAMBLE
\providecommand{\DIFdel}[1]{{\protect\color{red}\sout{#1}}}                      %DIF PREAMBLE
%DIF SAFE PREAMBLE %DIF PREAMBLE
\providecommand{\DIFaddbegin}{} %DIF PREAMBLE
\providecommand{\DIFaddend}{} %DIF PREAMBLE
\providecommand{\DIFdelbegin}{} %DIF PREAMBLE
\providecommand{\DIFdelend}{} %DIF PREAMBLE
\providecommand{\DIFmodbegin}{} %DIF PREAMBLE
\providecommand{\DIFmodend}{} %DIF PREAMBLE
%DIF FLOATSAFE PREAMBLE %DIF PREAMBLE
\providecommand{\DIFaddFL}[1]{\DIFadd{#1}} %DIF PREAMBLE
\providecommand{\DIFdelFL}[1]{\DIFdel{#1}} %DIF PREAMBLE
\providecommand{\DIFaddbeginFL}{} %DIF PREAMBLE
\providecommand{\DIFaddendFL}{} %DIF PREAMBLE
\providecommand{\DIFdelbeginFL}{} %DIF PREAMBLE
\providecommand{\DIFdelendFL}{} %DIF PREAMBLE
%DIF COLORLISTINGS PREAMBLE %DIF PREAMBLE
\RequirePackage{listings} %DIF PREAMBLE
\RequirePackage{color} %DIF PREAMBLE
\lstdefinelanguage{DIFcode}{ %DIF PREAMBLE
%DIF DIFCODE_UNDERLINE %DIF PREAMBLE
  moredelim=[il][\color{red}\sout]{\%DIF\ <\ }, %DIF PREAMBLE
  moredelim=[il][\color{blue}\uwave]{\%DIF\ >\ } %DIF PREAMBLE
} %DIF PREAMBLE
\lstdefinestyle{DIFverbatimstyle}{ %DIF PREAMBLE
	language=DIFcode, %DIF PREAMBLE
	basicstyle=\ttfamily, %DIF PREAMBLE
	columns=fullflexible, %DIF PREAMBLE
	keepspaces=true %DIF PREAMBLE
} %DIF PREAMBLE
\lstnewenvironment{DIFverbatim}{\lstset{style=DIFverbatimstyle}}{} %DIF PREAMBLE
\lstnewenvironment{DIFverbatim*}{\lstset{style=DIFverbatimstyle,showspaces=true}}{} %DIF PREAMBLE
%DIF END PREAMBLE EXTENSION ADDED BY LATEXDIFF

\begin{document}

\maketitle
\DIFdelbegin %DIFDELCMD < 

%DIFDELCMD < %%%
\DIFdelend \section*{Outline and motivation}

We now turn our attention to the generation of seismic waves during an \textbf{earthquake}. In doing so, we see 
immediately that our equations of motion imply that earthquakes cannot happen. \DIFdelbegin \DIFdel{To rectify things, 
we will not produce }\DIFdelend \DIFaddbegin \DIFadd{We will not attempt to rectify this by 
producing }\DIFaddend a complete \textbf{dynamical theory} for earthquakes\DIFaddbegin \DIFadd{, }\DIFaddend as no such theory yet exists. Instead, we 
\DIFdelbegin \DIFdel{just seek to }\DIFdelend \DIFaddbegin \DIFadd{simply }\DIFaddend modify our equations in a \textbf{phenomenological} manner to account for a specified displacement 
across a \textbf{fault}. Because this displacement is regarded as a parameter and not as a dynamic variable, 
this approach is known as the \textbf{kinematic source approximation}. The resulting theory has been found to work very 
well in practice, and it allows us to set up an \textbf{inverse problem} to determine the displacement across the 
fault from recorded seismic waves. In the final part of the lecture we begin to think more about inverse 
problems, this being a topic we will return to throughout these lectures.

\section*{The elastic rebound theory for earthquakes}

\begin{figure}[ht]
    \centering     
     \DIFdelbeginFL %DIFDELCMD < \includegraphics[width=0.7\textwidth]{figures/L3F1.png} 
%DIFDELCMD <     %%%
\DIFdelendFL \DIFaddbeginFL \includegraphics[width=0.7\textwidth]{figures/L3F1.pdf} 
    \DIFaddendFL \caption{A schematic \DIFdelbeginFL \DIFdelFL{diagram of }\DIFdelendFL \DIFaddbeginFL \DIFaddFL{cross-section through }\DIFaddendFL a \DIFaddbeginFL \DIFaddFL{normal }\DIFaddendFL fault. You do not need to learn the associated terminology.}
    \label{fig:1}
\end{figure}

The majority of natural earthquakes occur due to motion across \textbf{faults} within the Earth \DIFaddbegin \DIFadd{(see Fig.~\ref{fig:1})}\DIFaddend , with this motion resulting from the
slow \DIFdelbegin \DIFdel{build up}\DIFdelend \DIFaddbegin \DIFadd{build-up}\DIFaddend , and then rapid release, of energy accumulated during long-term tectonic deformation. This long-term
deformation is associated with processes such as mantle convection and plate tectonics. Indeed, most earthquakes occur at or 
near plate boundaries, and the type of motion across the fault is determined broadly by the plate tectonic setting. 
The exceptions to this picture are (i) earthquakes associated with motion of fluids within volcanic systems and (ii) deep earthquakes 
occurring within subducted slabs whose cause remains uncertain\DIFdelbegin \footnote{\DIFdel{At these depths it is thought that 
rocks will be too hot to undergo brittle failure, and so the process must seemingly be different than for shallow earthquakes. 
In terms of the seismic waves generated, however, there is little to distinguish deep earthquakes from shallower 
ones caused by brittle failure.}}%DIFAUXCMD
\addtocounter{footnote}{-1}%DIFAUXCMD
\DIFdel{.}\DIFdelend \DIFaddbegin \DIFadd{.}\footnote{\DIFadd{At these depths it is thought that 
rocks will be too hot to undergo brittle failure, and so the process must seemingly be different from that for shallow earthquakes. 
In terms of the seismic waves generated, however, there is little to distinguish deep earthquakes from shallower 
ones caused by brittle failure.}} \DIFaddend There are also man-made earthquakes, including nuclear weapon tests, \DIFdelbegin \DIFdel{and }\DIFdelend explosions generated 
as part of oil prospecting, and small \DIFdelbegin \footnote{\DIFdel{But sometimes still quite damaging.}} %DIFAUXCMD
\addtocounter{footnote}{-1}%DIFAUXCMD
\DIFdelend earthquakes associated with fracking or extraction of natural gas.\DIFaddbegin \footnote{\DIFadd{These are small, but sometimes still quite damaging.}}
\DIFaddend 

The idea that earthquakes are produced by rapid motion across a fault was formulated by Harry Fielding Reid (\DIFdelbegin \DIFdel{1859-1944}\DIFdelend \DIFaddbegin \DIFadd{1859--1944}\DIFaddend ) following
a large earthquake in San Francisco in 1906. Based on 
observations around San Francisco, Reid arrived at his \textbf{elastic rebound theory} \DIFaddbegin \DIFadd{(Reid 1910) }\DIFaddend as summarised in \DIFdelbegin \DIFdel{fig.2}\DIFdelend \DIFaddbegin \DIFadd{Fig.~\ref{fig:2}}\DIFaddend . Prior 
to the earthquake, deformation is gradually built up on either 
side of the locked fault due to tectonic processes, and this stores energy within the system. At some point, the tractions across 
the fault become sufficiently large to overcome 
frictional resistance, rupture occurs rapidly, and seismic waves radiate away some of the stored energy.

While this \DIFdelbegin \DIFdel{"theory" }\DIFdelend \DIFaddbegin \DIFadd{``theory'' }\DIFaddend is light on details, it is surely correct in a broad sense, 
and subsequent work has tried to make things more precise. Key questions include: (i) the nature of frictional resistance on 
the fault, (ii) the conditions under which rupture will occur,
and (iii) the factors that determine when rupture stops. Why, for example, are some earthquakes on the same fault small and others big? 
Such work also informs ideas about earthquake predictability, though \DIFdelbegin \DIFdel{the consensus in the field }\DIFdelend \DIFaddbegin \DIFadd{a widely held view }\DIFaddend is that earthquakes are 
a so-called \textbf{critical \DIFdelbegin \DIFdel{phenomena}\DIFdelend \DIFaddbegin \DIFadd{phenomenon}\DIFaddend }, with the size of a particular earthquake being so sensitive to initial 
conditions that meaningful predictions are impossible. 
The physics of earthquakes is a large and difficult 
subject, about which much remains to be learned. In this course, however, we will approximate the full physics by regarding 
the motion on the fault as simply a parameter within the theory, and calculate the seismic waves generated. Within 
this \textbf{kinematic source approximation}, any interaction between the radiated waves and rupture propagation 
is neglected, but it is thought that this is a good approximation in most circumstances.

\begin{figure}[ht]
    \centering    
     \DIFdelbeginFL %DIFDELCMD < \includegraphics[width=0.7\textwidth]{figures/L3F2.png} 
%DIFDELCMD <     %%%
\DIFdelendFL \DIFaddbeginFL \includegraphics[width=\textwidth]{figures/L3F2.pdf} 
    \DIFaddendFL \caption{A graphical representation of Reid's elastic rebound theory.}
    \label{fig:2}
\end{figure}

\section*{Introduction of the stress glut}

Our aim is to incorporate motion on a fault into the equations of motion for an elastic body. To do this, we will 
actually take a more general, and yet ultimately simpler, approach. \DIFdelbegin \DIFdel{Returning to the material of Lecture12, }\DIFdelend \DIFaddbegin \DIFadd{In Lecture~13 }\DIFaddend we introduced the Lagrangian density
\DIFdelbegin \begin{displaymath}
\DIFdel{\mathcal{L}(\mathbf{x}, t, \boldsymbol{\varphi}, \mathbf{v}, \mathbf{F}) = \frac{1}{2}\rho(\mathbf{x})||\mathbf{v}||^2 - W(\mathbf{x}, \mathbf{F}),
}\end{displaymath}%DIFAUXCMD
\DIFdelend \DIFaddbegin \begin{equation}\DIFadd{\label{eq:1}
\mathcal{L}(\mathbf{x}, t, \bphi, \mathbf{v}, \mathbf{F}) = \frac{1}{2}\rho(\mathbf{x})\|\mathbf{v}\|^2 - W(\mathbf{x}, \mathbf{F}),
}\end{equation}\DIFaddend 
for the finite motion of a hyperelastic body, while the associated action is
\DIFdelbegin \begin{displaymath}
\DIFdel{\mathcal{S}(\boldsymbol{\varphi}) = \int_{0}^{T}\int_{M}\mathcal{L}[\mathbf{x}, t, \boldsymbol{\varphi}(\mathbf{x}, t), \mathbf{v}(\mathbf{x}, t), \mathbf{F}(\mathbf{x}, t)]\dd^{3}\mathbf{x}\dd t,
}\end{displaymath}%DIFAUXCMD
\DIFdelend \DIFaddbegin \begin{equation}\DIFadd{\label{eq:2}
\mathcal{S}(\bphi) = \int_{0}^{T}\int_{M}\mathcal{L}[\mathbf{x}, t, \bphi(\mathbf{x}, t), \mathbf{v}(\mathbf{x}, t), \mathbf{F}(\mathbf{x}, t)]\dd^{3}\mathbf{x}\dd t,
}\end{equation}\DIFaddend 
where $[0, T]$ is the \DIFdelbegin \DIFdel{time-interval }\DIFdelend \DIFaddbegin \DIFadd{time interval }\DIFaddend of interest. From the resulting equations of motion it \DIFdelbegin \DIFdel{is readily shown }\DIFdelend \DIFaddbegin \DIFadd{can be shown (and you will do so in the first problem set) }\DIFaddend that the 
\textbf{total energy} of the elastic body is conserved\DIFdelbegin \begin{displaymath}
\DIFdel{\frac{\ddns}{\ddns t}\int_{M}E(\mathbf{x}, t)\dd^{3}\mathbf{x} = 0,
}\end{displaymath}%DIFAUXCMD
\DIFdelend \DIFaddbegin \DIFadd{,
}\begin{equation}\DIFadd{\label{eq:3}
\frac{\ddns}{\ddns t}\int_{M}E(\mathbf{x}, t)\dd^{3}\mathbf{x} = 0,
}\end{equation}\DIFaddend 
where the \textbf{energy density} is defined in the expected manner\DIFdelbegin \begin{displaymath}
\DIFdel{E(\mathbf{x}, t) = \frac{1}{2}\rho(\mathbf{x})||\mathbf{v}(\mathbf{x}, t)||^2 + W}[\DIFdel{\mathbf{x}, \mathbf{F}(\mathbf{x}, t)}]\DIFdel{.
}\end{displaymath}%DIFAUXCMD
\DIFdelend \DIFaddbegin \DIFadd{,
}\begin{equation}\DIFadd{\label{eq:4}
E(\mathbf{x}, t) = \frac{1}{2}\rho(\mathbf{x})\|\mathbf{v}(\mathbf{x}, t)\|^2 + W}[\DIFadd{\mathbf{x}, \mathbf{F}(\mathbf{x}, t)}]\DIFadd{.
}\end{equation}\DIFaddend 
If the body is initially in a stable equilibrium configuration, then any spontaneous non-trivial \DIFdelbegin \footnote{\DIFdel{Meaning, not a rigid-body motion, 
but one that causes a change in shape and/or volume.}} %DIFAUXCMD
\addtocounter{footnote}{-1}%DIFAUXCMD
\DIFdelend deformation will, by definition, raise the potential energy\DIFdelbegin \begin{displaymath}
\DIFdel{\mathcal{V}(t) = \int_{M}W}[\DIFdel{\mathbf{x}, \mathbf{F}(\mathbf{x}, t)}]\DIFdel{\dd^{3}\mathbf{x}.
}\end{displaymath}%DIFAUXCMD
\DIFdelend \DIFaddbegin \DIFadd{,}\footnote{\DIFadd{By non-trivial we mean a deformation that is not a rigid-body motion, 
but one that causes a change in shape and/or volume.}}
\begin{equation}\DIFadd{\label{eq:5}
\mathcal{V}(t) = \int_{M}W}[\DIFadd{\mathbf{x}, \mathbf{F}(\mathbf{x}, t)}]\DIFadd{\dd^{3}\mathbf{x}.
}\end{equation}\DIFaddend 
Because the kinetic energy is non-negative, this must lead to the overall energy of the body being increased. Eq.\DIFdelbegin \DIFdel{(3}\DIFdelend \DIFaddbegin \DIFadd{~(\ref{eq:3}}\DIFaddend ) implies that such an 
increase in the total energy cannot occur, and so if a body is initially at rest, then it will remain so for all
later times. 

It is clear that our theory needs to be modified in some way to account for earthquakes. The potential energy is the obvious place to look 
as this is the part whose form was essentially guessed at. We recall that the principle of material frame indifference requires that 
the potential energy $\mathcal{V}(t)$ of the body be:
\begin{enumerate}
    \item invariant under time-translations;
    \item invariant under superimposed rigid body translations;
    \item invariant under superimposed rigid body rotations.
\end{enumerate}
A kinematic description of a seismic source necessarily breaks the first of these invariance properties, as the onset of motion on 
the fault gives a \textbf{distinguished point in time}. We shall still retain the final two invariance properties of the potential 
energy. Indeed, they are necessary for the seismic source to represent an internal process that does not alter the net linear \DIFdelbegin \DIFdel{nor 
}\DIFdelend \DIFaddbegin \DIFadd{or 
}\DIFaddend angular momentum of the Earth. Having dropped the time-translation invariance, we can generalise the Lagrangian density to become
\DIFdelbegin \begin{displaymath}
\DIFdel{\mathcal{L}(\mathbf{x}, t, \boldsymbol{\varphi}, \mathbf{v}, \mathbf{F}) = \frac{1}{2}\rho(\mathbf{x})||\mathbf{v}||^2 - W(\mathbf{x}, t, \mathbf{F}),
}\end{displaymath}%DIFAUXCMD
\DIFdelend \DIFaddbegin \begin{equation}\DIFadd{\label{eq:6}
\mathcal{L}(\mathbf{x}, t, \bphi, \mathbf{v}, \mathbf{F}) = \frac{1}{2}\rho(\mathbf{x})\|\mathbf{v}\|^2 - W(\mathbf{x}, t, \mathbf{F}),
}\end{equation}\DIFaddend 
where we note that $W(\mathbf{x}, t, \mathbf{F}) = W(\mathbf{x}, t, \mathbf{Q}\mathbf{F})$ for all rotation matrices $\mathbf{Q}$. It is 
clear that the resulting \DIFdelbegin \DIFdel{Euler-Lagrange equations remains unchanged}\begin{displaymath}
\DIFdel{\rho\frac{\partial v_{i}}{\partial t} - \frac{\partial T_{ij}}{\partial x_{j}} = 0
}\end{displaymath}%DIFAUXCMD
\DIFdelend \DIFaddbegin \DIFadd{Euler--Lagrange equations remain unchanged,
}\begin{equation}\DIFadd{\label{eq:7}
\rho\frac{\partial v_{i}}{\partial t} - \frac{\partial T_{ij}}{\partial x_{j}} = 0,
}\end{equation}\DIFaddend 
and similarly for the boundary conditions
\DIFdelbegin \begin{displaymath}
\DIFdel{T_{ij}\hat{n}_{j} = 0
}\end{displaymath}%DIFAUXCMD
\DIFdelend \DIFaddbegin \begin{equation}\DIFadd{\label{eq:8}
T_{ij}\hat{n}_{j} = 0
}\end{equation}\DIFaddend 
on $\partial M$. However, the first \DIFdelbegin \DIFdel{Piola-Kirchhoff }\DIFdelend \DIFaddbegin \DIFadd{Piola--Kirchhoff }\DIFaddend stress tensor now depends on the time-dependent strain energy function through the 
generalised constitutive relation
\DIFdelbegin \begin{displaymath}
\DIFdel{T_{ij}(\mathbf{x}, t) = \frac{\partial W}{\partial F_{ij}}[\mathbf{x}, t, \mathbf{F}(\mathbf{x}, t)].
}\end{displaymath}%DIFAUXCMD
\DIFdelend \DIFaddbegin \begin{equation}\DIFadd{\label{eq:9}
T_{ij}(\mathbf{x}, t) = \frac{\partial W}{\partial F_{ij}}[\mathbf{x}, t, \mathbf{F}(\mathbf{x}, t)].
}\end{equation}\DIFaddend 
Within Lecture\DIFdelbegin \DIFdel{13}\DIFdelend \DIFaddbegin \DIFadd{~14}\DIFaddend , we saw that the invariance of the potential energy under superimposed rigid body rotations requires that the strain 
energy function depends on the deformation gradient only through the right 
\DIFdelbegin \DIFdel{Cauchy-Green }\DIFdelend \DIFaddbegin \DIFadd{Cauchy--Green }\DIFaddend deformation tensor. Using this result, we shall assume that the time-dependent strain energy function can be decomposed 
into a sum of two terms\DIFdelbegin \begin{displaymath}
\DIFdel{W(\mathbf{x}, t, \mathbf{F}) = U(\mathbf{x}, \mathbf{C}) - \frac{1}{2}\overline{S}_{ij}(\mathbf{x}, t)C_{ij}
}\end{displaymath}%DIFAUXCMD
\DIFdelend \DIFaddbegin \DIFadd{,
}\begin{equation}\DIFadd{\label{eq:10}
W(\mathbf{x}, t, \mathbf{F}) = U(\mathbf{x}, \mathbf{C}) - \frac{1}{2}\overline{S}_{ij}(\mathbf{x}, t)C_{ij},
}\end{equation}\DIFaddend 
where $\overline{S}_{ij}$ is a given time-dependent symmetric second-order tensor field. Here $U$ describes the assumed constitutive 
properties of the material, which do not change with time, while $\overline{S}_{ij}$ represents the effects of the seismic source. 
Though this is not the most general possible modification to the strain energy,
more general forms do not change the linearised equations of motion. Using eq.\DIFdelbegin \DIFdel{(9}\DIFdelend \DIFaddbegin \DIFadd{~(\ref{eq:9}}\DIFaddend ), we see that eq.\DIFdelbegin \DIFdel{(10) implies
}\begin{displaymath}
\DIFdel{T_{ij}(\mathbf{x}, t) = F_{ik}(\mathbf{x}, t)\left\{2\frac{\partial U}{\partial C_{kj}}[\mathbf{x}, \mathbf{C}(\mathbf{x}, t)] - 
\overline{S}_{kj}(\mathbf{x}, t)\right\}.
}\end{displaymath}%DIFAUXCMD
\DIFdelend \DIFaddbegin \DIFadd{~(\ref{eq:10}) implies
}\begin{equation}\DIFadd{\label{eq:11}
T_{ij}(\mathbf{x}, t) = F_{ik}(\mathbf{x}, t)\left\{2\frac{\partial U}{\partial C_{kj}}[\mathbf{x}, \mathbf{C}(\mathbf{x}, t)] - 
\overline{S}_{kj}(\mathbf{x}, t)\right\},
}\end{equation}\DIFaddend 
and \DIFaddbegin \DIFadd{we }\DIFaddend can use this relation to introduce the \textbf{stress glut}, $\overline{T}_{ij}$, through the identity
\DIFdelbegin \begin{displaymath}
\DIFdel{\overline{T}_{ij}(\mathbf{x}, t) = F_{ik}(\mathbf{x}, t)\overline{S}_{kj}(\mathbf{x}, t) = 2F_{ik}(\mathbf{x}, t)
\frac{\partial U}{\partial C_{kj}}[\mathbf{x}, \mathbf{C}(\mathbf{x}, t)] - T_{ij}(\mathbf{x}, t).
}\end{displaymath}%DIFAUXCMD
\DIFdelend \DIFaddbegin \begin{equation}\DIFadd{\label{eq:12}
\overline{T}_{ij}(\mathbf{x}, t) = F_{ik}(\mathbf{x}, t)\overline{S}_{kj}(\mathbf{x}, t) = 2F_{ik}(\mathbf{x}, t)
\frac{\partial U}{\partial C_{kj}}[\mathbf{x}, \mathbf{C}(\mathbf{x}, t)] - T_{ij}(\mathbf{x}, t).
}\end{equation}\DIFaddend 
Physically, the stress glut represents the difference between the stress predicted by the elastic constitutive relation 
and its true value. It can, therefore, be interpreted as arising 
from a \textbf{failure of elasticity to correctly determine stresses within the body}. This point of view is consistent with 
elastic rebound theory, with material near to a fault undoubtedly failing to be elastic.

\section*{Linearised equations of motion incorporating a stress glut}

We consider a simple hyperelastic body that is initially at rest in a stable equilibrium configuration. At some 
later time, a non-zero stress glut is initiated, and we wish to model the subsequent linearised motion of the body. As 
with Lecture\DIFdelbegin \DIFdel{13}\DIFdelend \DIFaddbegin \DIFadd{~14}\DIFaddend , we assume that the motion can be written
\DIFdelbegin \begin{displaymath}
\DIFdel{\varphi_{i}(\mathbf{x}, t) = x_{i} + s u_{i}(\mathbf{x}, t) + O(s^2),
}\end{displaymath}%DIFAUXCMD
\DIFdelend \DIFaddbegin \begin{equation}\DIFadd{\label{eq:13}
\varphi_{i}(\mathbf{x}, t) = x_{i} + s u_{i}(\mathbf{x}, t) + O(s^2),
}\end{equation}\DIFaddend 
where we note that the reference body has been chosen such that the equilibrium configuration is the identity mapping. 
Without loss of generality, the initial conditions on the displacement vector can be applied at time $t=0$, and so we have
\DIFdelbegin \begin{displaymath}
\DIFdel{u_{i}(\mathbf{x}, 0) = 0, \quad \frac{\partial u_{i}}{\partial t}(\mathbf{x}, 0) = 0.
}\end{displaymath}%DIFAUXCMD
\DIFdelend \DIFaddbegin \begin{equation}\DIFadd{\label{eq:14}
u_{i}(\mathbf{x}, 0) = 0,\quad \frac{\partial u_{i}}{\partial t}(\mathbf{x}, 0) = 0.
}\end{equation}\DIFaddend 
Motivated by eq.\DIFdelbegin \DIFdel{(10}\DIFdelend \DIFaddbegin \DIFadd{~(\ref{eq:10}}\DIFaddend ), we write the strain energy function for the body in the form
\DIFdelbegin \begin{displaymath}
\DIFdel{W(\mathbf{x}, t, \mathbf{F}) = U(\mathbf{x}, \mathbf{C}) - s\frac{1}{2}\overline{S}_{ij}(\mathbf{x}, t)C_{ij},
}\end{displaymath}%DIFAUXCMD
\DIFdelend \DIFaddbegin \begin{equation}\DIFadd{\label{eq:15}
W(\mathbf{x}, t, \mathbf{F}) = U(\mathbf{x}, \mathbf{C}) - \frac{s}{2}\overline{S}_{ij}(\mathbf{x}, t)C_{ij},
}\end{equation}\DIFaddend 
where we note that the time-dependent term has no zeroth-order part. Following the method in Lecture\DIFdelbegin \DIFdel{13}\DIFdelend \DIFaddbegin \DIFadd{~14}\DIFaddend , it is a simple 
matter to obtain the corresponding linearised equations of motion\DIFdelbegin \begin{displaymath}
\DIFdel{\rho\frac{\partial^2 u_{i}}{\partial t^2} - \frac{\partial}{\partial x_{j}}\left(A_{ijkl}\frac{\partial u_{k}}{\partial x_{l}}\right) = -\frac{\partial\overline{S}_{ij}}{\partial x_{j}},
}\end{displaymath}%DIFAUXCMD
\DIFdelend \DIFaddbegin \DIFadd{. In doing so, the elastic tensor $A_{ijkl}$ is obtained as the second 
derivative with respect to $\mathbf{F}$, evaluated at $\mathbf{F}=\mathbf{I}$, of the time-independent part $U(\mathbf{x},\mathbf{C})$ of 
the strain energy function alone; the term involving the stress glut is already of first order in $s$, and contributes only to 
the force terms. This is the sense in which the definition of the elastic tensor given in Lecture~14 is to be understood 
whenever a seismic source is present. We find
}\begin{equation}\DIFadd{\label{eq:16}
\rho\frac{\partial^2 u_{i}}{\partial t^2} - \frac{\partial}{\partial x_{j}}\left(A_{ijkl}\frac{\partial u_{k}}{\partial x_{l}}\right) = -\frac{\partial\overline{S}_{ij}}{\partial x_{j}},
}\end{equation}\DIFaddend 
and boundary conditions 
\DIFdelbegin \begin{displaymath}
\DIFdel{A_{ijkl}\frac{\partial u_{k}}{\partial x_{l}}\hat{n}_{j} = \overline{S}_{ij}\hat{n}_{j}
}\end{displaymath}%DIFAUXCMD
\DIFdelend \DIFaddbegin \begin{equation}\DIFadd{\label{eq:17}
A_{ijkl}\frac{\partial u_{k}}{\partial x_{l}}\hat{n}_{j} = \overline{S}_{ij}\hat{n}_{j}
}\end{equation}\DIFaddend 
on $\partial M$. In this manner, we see that the seismic source enters into the linearised equations of motion as an 
\textbf{effective body force and an effective surface traction}, with 
each determined from the stress glut. Note that within this linearised theory the tensors $\overline{S}_{ij}$ and 
$\overline{T}_{ij}$ agree to first order, with \textbf{either being called the stress glut}.
In particular, within the linearised equations of motion, the \textbf{stress glut is necessarily symmetric}.

\section*{The stress glut for an idealised fault}


\begin{figure}[ht]
    \centering    
     \DIFdelbeginFL %DIFDELCMD < \includegraphics[width=0.7\textwidth]{figures/L3F3.png} 
%DIFDELCMD <     %%%
\DIFdelendFL \DIFaddbeginFL \includegraphics[width=0.55\textwidth]{figures/L3F3.png} 
    \DIFaddendFL \caption{A fault in cross section, where the offset of the strata on either side can be clearly
seen. It is also apparent that the fault is not a single surface along which the
motion is discontinuous, but rather a zone of \DIFdelbeginFL \DIFdelFL{finite-width }\DIFdelendFL \DIFaddbeginFL \DIFaddFL{finite width }\DIFaddendFL which has undergone
severe localised deformation. \DIFaddbeginFL \DIFaddFL{Photograph: source unknown; the road at the bottom gives the scale, the exposure being a few metres high.}\DIFaddendFL }
    \label{fig:3}
\end{figure}

We wish to determine the appropriate stress glut for an \textbf{idealised fault}. On a real fault, there will be large 
non-elastic deformation within some \textbf{finite fault zone}\DIFaddbegin \DIFadd{, as can be seen in Fig.~\ref{fig:3}}\DIFaddend . We shall, 
however, neglect such complexities, and suppose that the elastic constitutive relation fails only on the fault 
surface across which the displacement is discontinuous. Such an approximation 
is appropriate so long as we consider waves having wavelength substantially larger \DIFdelbegin \DIFdel{that }\DIFdelend \DIFaddbegin \DIFadd{than }\DIFaddend the width of the fault zone, 
and this is true for most seismological applications (e.g.\DIFaddbegin \DIFadd{\ }\DIFaddend seismic wavelengths of kilometres and fault \DIFdelbegin \DIFdel{zones }\DIFdelend \DIFaddbegin \DIFadd{zone }\DIFaddend widths of metres).

We assume that the \textbf{fault plane}, $\Sigma$, passes through the origin, and lies within the $(x_1, x_2)$-plane. The discontinuity 
in the displacement vector on the fault plane is assumed to be only in its tangential component so that there is no separation or 
interpenetration of material, and hence \DIFdelbegin \DIFdel{it }\DIFdelend \DIFaddbegin \DIFadd{the displacement vector }\DIFaddend can be decomposed in the form
\DIFdelbegin \begin{displaymath}
\DIFdel{u_{i}(x_1, x_2, x_3, t) = u_{i}^{0}(x_1, x_2, x_3, t) + w_{i}(x_1, x_2, t)H(x_3),
}\end{displaymath}%DIFAUXCMD
\DIFdelend \DIFaddbegin \begin{equation}\DIFadd{\label{eq:18}
u_{i}(x_1, x_2, x_3, t) = u_{i}^{0}(x_1, x_2, x_3, t) + w_{i}(x_1, x_2, t)H(x_3),
}\end{equation}\DIFaddend 
where $u_i^0$ is continuous, $w_i$ is tangential on $\Sigma$ (i.e.\DIFaddbegin \DIFadd{\ }\DIFaddend $w_3=0$), and $H(\cdot)$ denotes the Heaviside step 
function. We call $\mathbf{w}$ the \textbf{displacement discontinuity vector}. 
We wish to determine the linearised first \DIFdelbegin \DIFdel{Piola-Kirchhoff }\DIFdelend \DIFaddbegin \DIFadd{Piola--Kirchhoff }\DIFaddend stress tensor associated with this displacement vector 
field using the assumed hyperelastic constitutive relation. To do this, 
we must first differentiate the displacement vector to obtain
\DIFdelbegin \begin{displaymath}
\DIFdel{\frac{\partial u_{i}}{\partial x_{j}} = \frac{\partial u_{i}^{0}}{\partial x_{j}} + \frac{\partial w_{i}}{\partial x_{j}}H(x_3) + w_{i}\delta_{j3}\delta(x_3),
}\end{displaymath}%DIFAUXCMD
\DIFdelend \DIFaddbegin \begin{equation}\DIFadd{\label{eq:19}
\frac{\partial u_{i}}{\partial x_{j}} = \frac{\partial u_{i}^{0}}{\partial x_{j}} + \frac{\partial w_{i}}{\partial x_{j}}H(x_3) + w_{i}\delta_{j3}\delta(x_3),
}\end{equation}\DIFaddend 
where we have used the well-known identity $H'(x) = \delta(x)$. It follows that the linearised stress tensor predicted by the constitutive relation takes the form
\DIFdelbegin \begin{displaymath}
\DIFdel{A_{ijkl}\frac{\partial u_{k}^{0}}{\partial x_{l}} + A_{ijkl}\frac{\partial w_{k}}{\partial x_{l}}H(x_3) + A_{ijk3}w_{k}\delta(x_3).
}\end{displaymath}%DIFAUXCMD
\DIFdelend \DIFaddbegin \begin{equation}\DIFadd{\label{eq:20}
A_{ijkl}\frac{\partial u_{k}^{0}}{\partial x_{l}} + A_{ijkl}\frac{\partial w_{k}}{\partial x_{l}}H(x_3) + A_{ijk3}w_{k}\delta(x_3).
}\end{equation}\DIFaddend 
The first term here is everywhere continuous, the second has a jump discontinuity \DIFaddbegin \DIFadd{across }\DIFaddend the fault plane, while the third is non-zero only 
on the fault plane but is singular there. 
That the stress might be discontinuous across the fault plane \DIFaddbegin \DIFadd{(in its tangential components; continuity of traction is enforced by 
the equations of motion themselves) }\DIFaddend is physically acceptable, with this property approximating rapid but 
continuous variations in stress across the \DIFdelbegin \DIFdel{finite-width }\DIFdelend \DIFaddbegin \DIFadd{finite width }\DIFaddend of a fault zone. It is, 
therefore, only the final singular term within eq.\DIFdelbegin \DIFdel{(20}\DIFdelend \DIFaddbegin \DIFadd{~(\ref{eq:20}}\DIFaddend ) that is non-physical and needs to be removed by inclusion of a stress 
glut. Specialising eq.\DIFdelbegin \DIFdel{(12}\DIFdelend \DIFaddbegin \DIFadd{~(\ref{eq:12}}\DIFaddend ) to the linearised case, we see that a suitable choice is given by
\DIFdelbegin \begin{displaymath}
\DIFdel{\overline{S}_{ij} = A_{ijk3}w_{k}\delta(x_3).
}\end{displaymath}%DIFAUXCMD
\DIFdelend \DIFaddbegin \begin{equation}\DIFadd{\label{eq:21}
\overline{S}_{ij} = A_{ijk3}w_{k}\delta(x_3).
}\end{equation}\DIFaddend 
It is worth emphasising that this result is really \textbf{more of a definition than something we have derived}. Indeed, an ideal 
fault is just a mathematical abstraction, providing a useful model for 
something far more complicated. These ideas extend immediately to a general fault whose surface, $\Sigma$, need not be planar. 
The appropriate stress glut is
\DIFdelbegin \begin{displaymath}
\DIFdel{\overline{S}_{ij} = A_{ijkl}w_{k}n_{l}\delta_{\Sigma},
}\end{displaymath}%DIFAUXCMD
\DIFdelend \DIFaddbegin \begin{equation}\DIFadd{\label{eq:22}
\overline{S}_{ij} = A_{ijkl}w_{k}\hat{n}_{l}\delta_{\Sigma},
}\end{equation}\DIFaddend 
where the displacement discontinuity vector is everywhere tangent to $\Sigma$, $\hat{\mathbf{n}}$ is the unit normal vector 
to the surface, and $\delta_{\Sigma}$ is the Dirac surface distribution for $\Sigma$ which is defined through the relation
\DIFdelbegin \begin{displaymath}
\DIFdel{\int_{M}f\,\delta_{\Sigma}\dd^{3}\mathbf{x} = \int_{\Sigma}f\,dS,
}\end{displaymath}%DIFAUXCMD
\DIFdelend \DIFaddbegin \begin{equation}\DIFadd{\label{eq:23}
\int_{M}f\,\delta_{\Sigma}\dd^{3}\mathbf{x} = \int_{\Sigma}f\dd S,
}\end{equation}\DIFaddend 
for any smooth $f$.

\section*{Formulating an inverse problem}
We have now assembled suitable equations of motion to describe seismic wave generation and propagation 
within the Earth. For convenience, we recall that the linearised equations of motion are
\DIFdelbegin \begin{displaymath}
\DIFdel{\rho\frac{\partial^2 u_{i}}{\partial t^2} - \frac{\partial}{\partial x_{j}}\left(A_{ijkl}\frac{\partial u_{k}}{\partial x_{l}}\right) = -\frac{\partial\overline{S}_{ij}}{\partial x_{j}},
}\end{displaymath}%DIFAUXCMD
\DIFdelend \DIFaddbegin \begin{equation}\DIFadd{\label{eq:24}
\rho\frac{\partial^2 u_{i}}{\partial t^2} - \frac{\partial}{\partial x_{j}}\left(A_{ijkl}\frac{\partial u_{k}}{\partial x_{l}}\right) = -\frac{\partial\overline{S}_{ij}}{\partial x_{j}},
}\end{equation}\DIFaddend 
along with boundary conditions
\DIFdelbegin \begin{displaymath}
\DIFdel{A_{ijkl}\frac{\partial u_{k}}{\partial x_{l}}\hat{n}_{j} = \overline{S}_{ij}\hat{n}_{j}
}\end{displaymath}%DIFAUXCMD
\DIFdelend \DIFaddbegin \begin{equation}\DIFadd{\label{eq:25}
A_{ijkl}\frac{\partial u_{k}}{\partial x_{l}}\hat{n}_{j} = \overline{S}_{ij}\hat{n}_{j},
}\end{equation}\DIFaddend 
while the stress glut for an ideal fault takes the form
\DIFdelbegin \begin{displaymath}
\DIFdel{\overline{S}_{ij} = A_{ijkl}w_{k}n_{l}\delta_{\Sigma}.
}\end{displaymath}%DIFAUXCMD
\DIFdelend \DIFaddbegin \begin{equation}\DIFadd{\label{eq:26}
\overline{S}_{ij} = A_{ijkl}w_{k}\hat{n}_{l}\delta_{\Sigma}.
}\end{equation}\DIFaddend 
Within these equations we can identify a number of \textbf{model parameters}:
\begin{enumerate}
    \item The referential density, $\rho$.
    \item The elastic tensor, $A_{ijkl}$.
    \item The displacement discontinuity vector, $w_i$.
    \item Geometric parameters describing the shape of the earth model\DIFaddbegin \DIFadd{, }\DIFaddend including any internal discontinuities\DIFdelbegin \DIFdel{and for the fault plane}\DIFdelend \DIFaddbegin \DIFadd{, and the geometry of the fault surface}\DIFaddend .
\end{enumerate}
Given all this information it is possible, with sufficient effort, to solve the equations\DIFaddbegin \DIFadd{, }\DIFaddend though this generally requires use of a computer. 

Suppose we have made seismic observations following an earthquake. Let \DIFdelbegin \DIFdel{$\mathbf{x}_i$ for $i=1, ..., n$ }\DIFdelend \DIFaddbegin \DIFadd{$\mathbf{x}^{(k)}$, $k=1,\dots,n$, }\DIFaddend denote the surface locations
of the seismometers, with each instrument providing recordings over some time interval \DIFdelbegin \DIFdel{, $I = [t_1, t_2]$}\DIFdelend \DIFaddbegin \DIFadd{$[t_1, t_2]$}\DIFaddend . As noted above, given 
suitable model parameters we can simulate seismic wave propagation to produce \textbf{synthetic seismograms} at each seismometer location. 
This is an example of solving a \textbf{forward problem}, which is to say performing a mapping from model parameters to observable quantities. 
To each forward problem there is an associated \textbf{inverse problem} in which we try to work backwards from the observations to the 
underlying model parameters.

Both the forward and inverse problems are concerned with transforming one kind of information into another. There is, however, a 
fundamental difference between these two problems, and this is best expressed through the concept of being well- or ill-posed. For a 
problem to be \textbf{well-posed}
the following three things must hold:
\begin{enumerate}
    \item For any given set of inputs, there exists an output.
    \item There is a unique output for each given input.
    \item The output depends continuously on the input.
\end{enumerate}
Note that these conditions only make sense sequentially; if the first does not hold, then the second \DIFdelbegin \DIFdel{does not make sense, etc}\DIFdelend \DIFaddbegin \DIFadd{is meaningless, and so on}\DIFaddend . A 
problem that is not well-posed is said to be \textbf{ill-posed}. 

These definitions can be illustrated by considering a simple example.
Suppose for distinct points $x_i$, \DIFdelbegin \DIFdel{$i=1, ..., n$ }\DIFdelend \DIFaddbegin \DIFadd{$i=1,\dots,n$ }\DIFaddend on the real line we are given corresponding function values $\{y_i\}_{i=1}^n$. 
We can then ask if there is a polynomial, $p(x)$\DIFaddbegin \DIFadd{, }\DIFaddend of degree $m$ such that
\DIFdelbegin \begin{displaymath}
\DIFdel{p(x_i) = y_i.
}\end{displaymath}%DIFAUXCMD
\DIFdelend \DIFaddbegin \begin{equation}\DIFadd{\label{eq:27}
p(x_i) = y_i.
}\end{equation}\DIFaddend 
Here the inputs to the problem are the values $\{y_i\}_{i=1}^n$, and the desired outputs are the $m+1$ coefficients of 
the polynomial. Considering first the existence question, it should be clear that if $m+1 \ge n$, then a suitable polynomial can always be 
found. If $m+1 < n$, then, in general, there is no such polynomial. But a solution is possible for certain inputs (e.g.\DIFaddbegin \DIFadd{\ }\DIFaddend the $y_i$ 
all lie on a line, and \DIFdelbegin \DIFdel{$m>1$}\DIFdelend \DIFaddbegin \DIFadd{$m\ge 1$}\DIFaddend ). Assuming existence has been established, we can turn to uniqueness. Here, again, we can 
readily see that the polynomial is unique if and only if \DIFdelbegin \DIFdel{$m+1 = n$}\DIFdelend \DIFaddbegin \DIFadd{$m+1 \le n$}\DIFaddend , while if $m+1 > n$ there are \DIFdelbegin \DIFdel{an infinite number of }\DIFdelend \DIFaddbegin \DIFadd{infinitely many }\DIFaddend possible answers. 
Finally, assuming that uniqueness has been established, it is possible to show that the coefficients of the polynomial 
depend continuously on the data, and hence a stable solution can be obtained.

A few additional comments can be made. First, reflecting the fact that all real measurements have uncertainties, we might choose 
to modify the problem to
\DIFdelbegin \begin{displaymath}
\DIFdel{p(x_i) = y_i + z_i,
}\end{displaymath}%DIFAUXCMD
\DIFdelend \DIFaddbegin \begin{equation}\DIFadd{\label{eq:28}
p(x_i) = y_i + z_i,
}\end{equation}\DIFaddend 
where $z_i$ are \textbf{unknown random errors} drawn from a known distribution. In addressing the existence question we must then 
broaden what it means for a solution to be found, with \DIFaddbegin \DIFadd{the }\DIFaddend fit to the data \DIFdelbegin \DIFdel{need not }\DIFdelend \DIFaddbegin \DIFadd{no longer required to }\DIFaddend be exact, but \DIFaddbegin \DIFadd{only adequate }\DIFaddend in a statistical sense.
As a result, our statements on existence in terms of $n$ and $m$ might need to be modified. Next, we can see that one potential 
cause of non-uniqueness is if the problem is \textbf{underdetermined}: we have fewer knowns than unknowns. For 
geophysical inverse problems this is almost always the case. As you might already imagine, uncertainty quantification within 
underdetermined problems is challenging, and this remains a major research topic.

Any plausible physical theory should be well-posed. If this were not the case, then it would not be possible to test it experimentally. 
For example, suppose that a physical theory satisfied the first two conditions above, 
but not the third. We could formulate an experiment to test our theory by comparing its predictions against observations made from 
known inputs. But the issue is that we can never know our inputs with complete precision. 
The lack of continuity within the theory means that tiny errors within these inputs can map onto 
arbitrarily large differences in the output, and this renders comparison between predictions and observations meaningless.

It can be shown that the seismological forward problem defined above is well-posed so long as the model parameters are subject to 
reasonable physical constraints (e.g.\DIFaddbegin \DIFadd{\ }\DIFaddend density should always be positive). By contrast, \textbf{the inverse problem is 
almost always ill-posed}. As noted above, the main reason is a lack of information. Real data sets are necessarily finite-dimensional,
but the model parameters are functions, and hence live in infinite-dimensional spaces. Later in this course 
we will see a number of practical methods for dealing with ill-posed inverse problems. None are perfect, and 
all must be applied with care. But these techniques do seem to have produced useful and reliable information about the Earth.

\section*{What you need to know and be able to do}
\DIFdelbegin %DIFDELCMD < \begin{itemize}
%DIFDELCMD <     %%%
\DIFdelend \DIFaddbegin \begin{enumerate}
    \DIFaddend \item[(i)] Know that most natural earthquakes are generated by rapid movement across faults within the Earth.
    \item[(ii)] Summarise Reid's elastic rebound theory for \DIFdelbegin \DIFdel{Earthquakes}\DIFdelend \DIFaddbegin \DIFadd{earthquakes}\DIFaddend .
    \item[(iii)] Explain how a seismic source can be introduced into elasticity using a stress glut, and outline the derivation of the stress glut for an idealised fault.
    \item[(iv)] Understand what is meant by an inverse problem along with the associated terminology and the definition of well- and ill-posed problems.
\DIFdelbegin %DIFDELCMD < \end{itemize}
%DIFDELCMD < %%%
\DIFdelend \DIFaddbegin \end{enumerate}
\DIFaddend 

\end{document}
